\documentclass[10pt,a4paper]{article}
\usepackage[margin=1.8cm]{geometry}
\usepackage{preamble}
\usepackage{amsmath,amssymb,amsthm}
\usepackage{multicol}
\usepackage{titlesec}
\usepackage{xcolor}
\usepackage{mdframed}
\usepackage{enumitem}
\usepackage{parskip}
\usepackage{booktabs}

% ── Colour palette ───────────────────────────────────────────────────────────
\definecolor{headbg}{HTML}{1B3A5C}
\definecolor{headfg}{HTML}{FFFFFF}
\definecolor{sectionbg}{HTML}{2E6DA4}
\definecolor{exbg}{HTML}{EBF4FB}
\definecolor{exframe}{HTML}{2E6DA4}
\definecolor{defbg}{HTML}{F5F5F5}
\definecolor{defframe}{HTML}{999999}

% ── Section titles ───────────────────────────────────────────────────────────


% ── Boxes ────────────────────────────────────────────────────────────────────
\newmdenv[backgroundcolor=defbg,linecolor=defframe,linewidth=0.8pt,
  innerleftmargin=6pt,innerrightmargin=6pt,innertopmargin=4pt,
  innerbottommargin=4pt,skipabove=3pt,skipbelow=3pt]{defbox}

\newmdenv[backgroundcolor=exbg,linecolor=exframe,linewidth=0.8pt,
  leftline=true,rightline=false,topline=false,bottomline=false,
  innerleftmargin=6pt,innerrightmargin=4pt,innertopmargin=3pt,
  innerbottommargin=3pt,skipabove=3pt,skipbelow=3pt]{exbox}

% ── Convenience macros ───────────────────────────────────────────────────────
\newcommand{\kw}[1]{\textbf{#1}}
\newcommand{\p}[2]{\frac{\partial #1}{\partial #2}}      % \p{f}{x}
\newcommand{\hvt}{\hat{\mathbf{t}}}
\newcommand{\hvn}{\hat{\mathbf{n}}}
\newcommand{\vx}{\mathbf{x}}
\newcommand{\vf}{\mathbf{f}}
\newcommand{\vg}{\mathbf{g}}
\newcommand{\eu}{\mathbf{e}_u}
\newcommand{\ev}{\mathbf{e}_v}
\newcommand{\ew}{\mathbf{e}_w}

\setlist[itemize]{leftmargin=*,topsep=1pt,itemsep=0pt,parsep=1pt}
\setlist[enumerate]{leftmargin=*,topsep=1pt,itemsep=0pt,parsep=1pt}

\setlength{\abovedisplayskip}{3pt}
\setlength{\belowdisplayskip}{3pt}
\setlength{\abovedisplayshortskip}{1pt}
\setlength{\belowdisplayshortskip}{1pt}

% ── Document ─────────────────────────────────────────────────────────────────
\begin{document}

\noindent\colorbox{headbg}{%
  \parbox{\linewidth}{%
    \vspace{6pt}
    \centering\color{headfg}\Large\bfseries
    MATH2811 Mathematical Methods II --- Part I Help Sheet\\[2pt]
    \normalsize
    Line Integrals $\cdot$ Surface \& Volume Integrals $\cdot$
    Differential Operators $\cdot$ Integral Theorems $\cdot$ Index Notation
    \vspace{6pt}
  }%
}

\vspace{4pt}
\begin{multicols}{2}

% ════════════════════════════════════════════════════════
\section{Line Integrals}

\begin{defbox}
The \kw{line integral of a scalar field} $f(\vx)$ along a curve $C$
with parametrisation $\vx(t)$, $t\in[t_0,t_1]$:
\[
  \int_C f\,d\ell
  = \int_{t_0}^{t_1} f(\vx(t))\left|\frac{d\vx}{dt}\right|dt.
\]
\textbf{Arclength:}\;
$L=\displaystyle\int_C d\ell
 = \int_{t_0}^{t_1}\!\left|\tfrac{d\vx}{dt}\right|dt$.
\quad
\textbf{Unit tangent:}\;
$\hvt=\dfrac{d\vx/dt}{\left|d\vx/dt\right|}$.

The \kw{line integral of a vector field} along an oriented curve:
\[
  \int_C \vf\cdot\hvt\,d\ell
  = \int_{t_0}^{t_1}\vf(\vx(t))\cdot\frac{d\vx}{dt}\,dt.
\]
For a \emph{closed} curve this is the \kw{circulation}
$\oint_C \vf\cdot\hvt\,d\ell$.
Both integrals are independent of parametrisation;
the vector integral changes sign if orientation is reversed.
\end{defbox}

% ════════════════════════════════════════════════════════
\section{Curvilinear Coordinates}

\begin{defbox}
\kw{Scale factors} and \kw{unit basis vectors} for orthogonal $(u,v,w)$:
\[
  h_u=\left|\p{\vx}{u}\right|,
  \quad
  \eu=\frac{1}{h_u}\p{\vx}{u},
  \quad\text{(similarly $v,w$).}
\]

\textbf{Cylindrical} $(r,\theta,z)$:
$h_r=1$, $h_\theta=r$, $h_z=1$.
\[
  \vx = r\cos\theta\,\mathbf{e}_1
      + r\sin\theta\,\mathbf{e}_2
      + z\,\mathbf{e}_3.
\]

\textbf{Spherical} $(r,\theta,\phi)$:
$h_r=1$, $h_\theta=r$, $h_\phi=r\sin\theta$.
\[
  \vx = r\sin\theta\cos\phi\,\mathbf{e}_1
      + r\sin\theta\sin\phi\,\mathbf{e}_2
      + r\cos\theta\,\mathbf{e}_3.
\]
\end{defbox}

% ════════════════════════════════════════════════════════
\section{Surface \& Volume Integrals}

\begin{defbox}
\kw{Unit normal} to surface $\vx(u,v)$:
$\;\hvn = \dfrac{\p{\vx}{u}\times\p{\vx}{v}}
               {\left|\p{\vx}{u}\times\p{\vx}{v}\right|}$.

\kw{Scalar surface integral} and \kw{surface area}:
\[
  \int_S f\,dS
  = \int_U f(\vx)\left|\p{\vx}{u}\times\p{\vx}{v}\right|du\,dv,
  \quad |S|=\int_S dS.
\]

\kw{Flux} of $\vf$ through oriented surface $S$:
\[
  \int_S\vf\cdot\hvn\,dS
  = \int_U\vf\cdot\!\left(\p{\vx}{u}\times\p{\vx}{v}\right)du\,dv.
\]

\kw{Volume integral}:
\[
  \int_V f\,dV
  = \int_U f\left|\p{\vx}{u}\cdot\!\left(\p{\vx}{v}\times\p{\vx}{w}\right)\right|du\,dv\,dw.
\]

In \textbf{orthogonal coords}:
$\bigl|\p{\vx}{u}\times\p{\vx}{v}\bigr|=h_uh_v$;\quad
$\bigl|\p{\vx}{u}\cdot(\p{\vx}{v}\times\p{\vx}{w})\bigr|=h_uh_vh_w$.

Jacobians: cylindrical $=r$;\quad spherical $=r^2\!\sin\theta$.
\end{defbox}

% ════════════════════════════════════════════════════════
\section{Differential Operators}

\begin{defbox}
\textbf{Gradient} (Cartesian):
\;$\bnabla f = \p{f}{x}\mathbf{e}_1+\p{f}{y}\mathbf{e}_2+\p{f}{z}\mathbf{e}_3$.

\textbf{Gradient} (orthogonal curvilinear):
\[
  \bnabla f
  = \frac{1}{h_u}\p{f}{u}\eu
  + \frac{1}{h_v}\p{f}{v}\ev
  + \frac{1}{h_w}\p{f}{w}\ew.
\]

\textbf{Divergence \& curl} (Cartesian):
\[
  \bnabla\cdot\vf
  = \p{f_1}{x}+\p{f_2}{y}+\p{f_3}{z},
\]
\[
  \bnabla\times\vf
  = \begin{vmatrix}
      \mathbf{e}_1 & \mathbf{e}_2 & \mathbf{e}_3\\
      \partial_x & \partial_y & \partial_z\\
      f_1 & f_2 & f_3
    \end{vmatrix}.
\]

\textbf{Divergence} (orthogonal curvilinear):
\begin{align*}
  \bnabla\cdot\vf &= \tfrac{1}{h_uh_vh_w}
     \bigl[\partial_u(h_vh_wf_u)\\
     &\quad+\partial_v(h_uh_wf_v)
     +\partial_w(h_uh_vf_w)\bigr].
\end{align*}
\textbf{Curl} (orthogonal curvilinear):
\[
  \bnabla\times\vf = \tfrac{1}{h_uh_vh_w}
  \begin{vmatrix}
    h_u\eu & h_v\ev & h_w\ew\\
    \partial_u & \partial_v & \partial_w\\
    h_uf_u & h_vf_v & h_wf_w
  \end{vmatrix}.
\]

\textbf{Key identities:}
$\bnabla\cdot(\bnabla\times\vf)\equiv 0$;\quad
$\bnabla\times(\bnabla f)\equiv\mathbf{0}$.

\textbf{Laplacian:}
$\bnabla^2g=\bnabla\cdot(\bnabla g)
 =\partial^2g/\partial x^2
  +\partial^2g/\partial y^2
  +\partial^2g/\partial z^2$.
\end{defbox}

\subsection{Product Rules}
\noindent($f,g$ scalar; $\vf,\vg$ vector; all $C^1$)
\begin{enumerate}[label=(\roman*)]
  \item $\bnabla(fg)\equiv g\bnabla f+f\bnabla g$
  \item $\bnabla\cdot(g\vf)\equiv(\bnabla g)\cdot\vf+g\,\bnabla\cdot\vf$
  \item $\bnabla\cdot(\vf\!\times\!\vg)\equiv\vg\cdot(\bnabla\!\times\!\vf)-\vf\cdot(\bnabla\!\times\!\vg)$
  \item $\bnabla\times(g\vf)\equiv g(\bnabla\!\times\!\vf)+(\bnabla g)\times\vf$
  \item $\bnabla\times(\vf\!\times\!\vg)\equiv(\vg\cdot\bnabla)\vf-(\vf\cdot\bnabla)\vg$\\
        $\hphantom{\bnabla\times(\vf\times\vg)\equiv{}}+\vf(\bnabla\cdot\vg)-\vg(\bnabla\cdot\vf)$
  \item $\bnabla(\vf\cdot\vg)\equiv(\vg\cdot\bnabla)\vf+(\vf\cdot\bnabla)\vg$\\
        $\hphantom{\bnabla(\vf\cdot\vg)\equiv{}}+\vf\times(\bnabla\!\times\!\vg)+\vg\times(\bnabla\!\times\!\vf)$
\end{enumerate}
\textbf{Vector Laplacian:}
$\bnabla\times(\bnabla\times\vf)=\bnabla(\bnabla\cdot\vf)-\bnabla^2\vf$.

% ════════════════════════════════════════════════════════
\section{Integral Theorems}

\begin{defbox}
\textbf{Fundamental Theorem of Line Integrals.}
For $C^1$ scalar $f$ and curve $\vx(t)$, $t\in[t_0,t_1]$:
\[
  \int_C \bnabla f\cdot\hvt\,d\ell
  = f(\vx(t_1))-f(\vx(t_0)).
\]
$\vf=\bnabla g$ is \kw{conservative}: path-independent line integrals;
on simply-connected domains, $\bnabla\times\vf=\mathbf{0}$.

\medskip
\textbf{Divergence Theorem (Gauss).}
$S$ closed surface, outward $\hvn$, bounding $V$:
\[
  \oint_S \vf\cdot\hvn\,dS = \int_V(\bnabla\cdot\vf)\,dV.
\]

\textbf{Stokes' Theorem.}
$C$ closed oriented curve bounding surface $S$:
\[
  \oint_C \vf\cdot\hvt\,d\ell
  = \int_S(\bnabla\times\vf)\cdot\hvn\,dS,
\]
where $\hvn$ is chosen so that $\hvt\times\hvn$ on $C$ is outward.
\end{defbox}

\begin{exbox}
\par\textbf{\color{exframe}Note.\space}
\textbf{Green's Theorem} (2D special case of Stokes'):
\[
  \oint_C \vf\cdot\hvt\,d\ell
  = \iint_A\!\left(\p{f_2}{x}-\p{f_1}{y}\right)dx\,dy,
\]
$C$ anti-clockwise, bounding $A\subset\R^2$.
\end{exbox}

% ════════════════════════════════════════════════════════
\section{Index Notation}

\begin{defbox}
Repeated index $\Rightarrow$ sum (\kw{Einstein convention}):
$\mathbf{a}\cdot\mathbf{b}=a_jb_j$.

\kw{Kronecker delta}: $\delta_{ij}=1$ if $i=j$, else $0$;
$\;a_i\delta_{ij}=a_j$;\;
$\delta_{ij}\delta_{jk}=\delta_{ik}$;\;
$\delta_{ii}=n$.

\kw{Levi-Civita}: $\varepsilon_{ijk}=+1$ (even perm.), $-1$ (odd), $0$ (repeat).
$\;[\mathbf{a}\times\mathbf{b}]_i=\varepsilon_{ijk}a_jb_k$.

\kw{$\varepsilon$--$\delta$ identity}:
\[
  \varepsilon_{ijk}\varepsilon_{klm}
  = \delta_{il}\delta_{jm}-\delta_{im}\delta_{jl}.
\]
Proves BAC--CAB:
$\mathbf{a}\times(\mathbf{b}\times\mathbf{c})
 =(\mathbf{a}\cdot\mathbf{c})\mathbf{b}-(\mathbf{a}\cdot\mathbf{b})\mathbf{c}$.

\textbf{Operators in index notation:}
\[
  [\bnabla f]_i=\p{f}{x_i},
  \quad
  \bnabla\cdot\vf=\p{f_j}{x_j},
  \quad
  [\bnabla\times\vf]_i=\varepsilon_{ijk}\p{f_k}{x_j}.
\]
\textbf{Antisymmetry trick:}
$\varepsilon_{ijk}\,\partial^2\!f_k/\partial x_i\partial x_j=0$
(antisym.\ $\times$ sym.\ $= 0$)
$\Rightarrow$ proves $\bnabla\cdot(\bnabla\times\vf)=0$ and
$\bnabla\times(\bnabla f)=\mathbf{0}$.
\end{defbox}

\subsection{Quick Reference}

\begin{center}
\small
\begin{tabular}{@{}ll@{}}
\toprule
\textbf{Object} & \textbf{Key formula / fact}\\
\midrule
Scalar line integral  & $\int_C f\,d\ell=\int f(\vx)|d\vx/dt|\,dt$\\
Vector line integral  & $\int_C\vf\cdot\hvt\,d\ell=\int\vf\cdot(d\vx/dt)\,dt$\\
Scale factors         & $h_u=|\p{\vx}{u}|$,\; $\eu=(1/h_u)\p{\vx}{u}$\\
Area element          & $|\partial_u\vx\times\partial_v\vx|=h_uh_v$\\
Volume element        & $|\partial_u\vx\cdot(\partial_v\vx\times\partial_w\vx)|=h_uh_vh_w$\\
FTLI                  & $\int_C\bnabla f\cdot\hvt\,d\ell=f(\vx(t_1))-f(\vx(t_0))$\\
Divergence Thm        & $\oint_S\vf\cdot\hvn\,dS=\int_V\bnabla\cdot\vf\,dV$\\
Stokes' Thm           & $\oint_C\vf\cdot\hvt\,d\ell=\int_S(\bnabla\times\vf)\cdot\hvn\,dS$\\
$\varepsilon$--$\delta$ & $\varepsilon_{ijk}\varepsilon_{klm}=\delta_{il}\delta_{jm}-\delta_{im}\delta_{jl}$\\
\bottomrule
\end{tabular}
\end{center}

\end{multicols}

\vspace{2pt}
\noindent\colorbox{headbg}{\parbox{\linewidth}{\centering\color{headfg}\small\vspace{2pt}
Conservative $\Leftrightarrow$ path-independent $\Leftrightarrow$ $\bnabla\times\vf=\mathbf{0}$
(simply-connected domain). \quad
All three integral theorems generalise the Fundamental Theorem of Calculus.
\vspace{2pt}}}

\end{document}